\subsection{AQ-67-EXACT: the canonical radius and the exact labelled $(0,4)$ image}
\paragraph{Scope and normalization.}
Consider the geodesic-cut interpretation of the map without marked-point
disks in cote 67/077. Fix a partition of four labelled punctures into two
pairs and remember the separating geodesic. Its length is $\ell>0$.
Cutting produces two hyperbolic pairs of pants of boundary lengths
$(\ell,0,0)$, where zero denotes a cusp. Cap the geodesic boundary by the
unit disk using its normalized hyperbolic arclength, as prescribed in
67/074. On each resulting three-pointed sphere use the unique coordinate
with cap centre $0$ and the two labelled cusps $1,\infty$.
Let $r(\ell)>0$ be the conformal radius of the cap at zero in that
coordinate. A cap coordinate $w$ has $|w|=1$ at the seam and
$z=r(\ell)w+O(w^2)$ after a rotation of $w$.
The tensor coordinate $q$ of 67/075 then has modulus $r(\ell)^2$.

Put
\[
 x=\frac{\ell}{4\pi},\qquad
 \phi(x)=\operatorname{Im}\bigl(\log\Gamma(1+ix)
                       -\log\Gamma(\tfrac12+ix)\bigr),
\]
where the analytic logarithm of $\Gamma$ on the right half-plane is
real on the positive real axis. Then
\[
 \boxed{\quad r(\ell)=4\exp\!\left(\frac{\phi(x)-\pi/4}{x}\right).\quad}
                                                               \tag{AQ.67.1}
\]
The function $r:(0,\infty)\longrightarrow(0,4)$ is a real-analytic
diffeomorphism. With the orientation of the normalized twist angle
$\theta\in\mathbb R/2\pi\mathbb Z$ chosen so that its increase multiplies
the tangent tensor by $e^{i\theta}$, the capping map is
\[
 (\ell,\theta)\longmapsto r(\ell)^2e^{i\theta}.
                                                               \tag{AQ.67.2}
\]
It is a real-analytic diffeomorphism onto $0<|q|<16$.
Its real tangent map is invertible everywhere, and its composite with
the inclusion $\{0<|q|<16\}\hookrightarrow\mathbb C^\times$ induces an
isomorphism of fundamental groups. These are statements about this
labelled geodesic-cut sector. They do not assert complex holomorphicity
or algebraic \emph{\'etaleness} of the map in 67/077--082, nor do they
settle the full multi-node question.

\paragraph{Geometric inputs and source locators.}
The hyperbolic pair of pants with specified nonnegative boundary
lengths is unique, with zero lengths interpreted as cusps. Its
Fenchel--Nielsen length and twist parametrization is recalled in
Mondello (2011), \S2.2, p.~640. The canonical attachment here is the
infinite flat-cylinder grafting of Mondello, \S5.4,
Theorem 5.4, p.~660; a punctured flat cylinder is conformally the
punctured disk with the arclength parametrization used above.
For an explicit independent source for the projective-connection
calculation below, see F\i rat (2021), equations (3.11), (3.13),
(3.14a), (3.23), and (3.30)--(3.34), arXiv:2102.03936v2,
pp.~14--17 and 24--25. We specialize the two other ends to cusps
and determine the branch directly below; we do not infer the branch
from that paper's integer-valued numerical prescription (3.27).

\begin{proof}
Write $\lambda=\ell/(2\pi)$, so $x=\lambda/2$.
The developing map $D$ of the hyperbolic metric is a local isometry
to a half-plane. Near a geodesic of length $2\pi\lambda$, diagonalize
its hyperbolic monodromy, with fixed points $0,\infty$.
In the arclength cap coordinate its analytic continuation is, up
to a nonzero constant, $w^{i\lambda}$. Indeed, increasing
$\arg w$ by $2\pi$ multiplies this function by $e^{-2\pi\lambda}$,
and along $|w|=1$ its hyperbolic arclength is
$\lambda\,d\arg w$. The same expression supplies the analytic
continuation across the seam. Its Schwarzian derivative therefore
has a double pole at the cap centre with leading coefficient
$(1+\lambda^2)/2$. At a cusp its leading coefficient is $1/2$,
as follows from the local developing coordinate $\log(z-z_0)$.
There are no other poles. On the sphere with centres and cusps
$0,1,\infty$, the resulting Schwarzian is consequently
\[
 T(z)=\frac{1+\lambda^2}{2z^2}
       +\frac1{2(z-1)^2}
       -\frac{1+\lambda^2}{2z(z-1)}.                           \tag{AQ.67.3}
\]
To see uniqueness, any two such projective connections differ by
a quadratic differential with at most simple poles at three points;
this space is $H^0(\mathbb P^1,\mathcal O(-1))=0$.

Set $a=(1+i\lambda)/2$. Direct substitution in
$y''+\tfrac12T(z)y=0$ gives the two solutions
\[
 y_\pm(z)=z^{(1\pm i\lambda)/2}(1-z)^{1/2}
     {}_2F_1\!\left(\frac{1\pm i\lambda}{2},
                    \frac{1\pm i\lambda}{2};
                    1\pm i\lambda;z\right).
\]
The ratio, normalized by its leading term at zero, is
\[
 D(z)=\frac{y_+(z)}{y_-(z)}
     =z^{i\lambda}
       \frac{{}_2F_1(a,a;2a;z)}
            {{}_2F_1(\bar a,\bar a;2\bar a;z)}.                \tag{AQ.67.4}
\]
All developing maps with the prescribed Schwarzian differ by a
M\"obius transformation. The ratio in (AQ.67.4) has the same
diagonal monodromy at zero; hence it is a valid developing map
after changing the target half-plane by that transformation.
Its boundary still contains the two fixed points $0,\infty$;
thus it is a straight half-plane bounded by a line through zero,
rather than a disk with a different boundary circle.

We now fix the phase without an unspecified multiple of $\pi$.
On the real interval $0<z<1$, $y_- =\overline{y_+}$.
Their Wronskian, computed at zero, is
$y_+'y_- -y_+y_-'=i\lambda$. In particular neither solution
vanishes there, and the continuously lifted argument $\Theta(z)$
of $D(z)$, normalized by
$\Theta(z)-\lambda\log z\longrightarrow0$ as $z\downarrow0$,
satisfies
\[
 \Theta'(z)=\frac{\lambda}{|y_+(z)|^2}>0.
                                                               \tag{AQ.67.5}
\]
The zero-balanced hypergeometric asymptotic
\[
 {}_2F_1(a,a;2a;z)
   \sim\frac{\Gamma(2a)}{\Gamma(a)^2}[-\log(1-z)]
       \qquad (z\uparrow1)
\]
is DLMF (15.4.21). Thus
\[
 \lim_{z\uparrow1}\Theta(z)=\vartheta(\lambda),\qquad
 \vartheta(\lambda)=2\operatorname{Im}
       \left(\log\Gamma(1+i\lambda)
                    -2\log\Gamma(\tfrac12+i\lambda/2)\right).
                                                               \tag{AQ.67.6}
\]
This is an equality of lifted arguments. For clarity, the
asymptotic is locally uniform in $\lambda\geq0$. The hypergeometric
factor on $0<z<1$ has no zeros when $\lambda>0$ by the Wronskian,
and at $\lambda=0$ it is positive. Its continuous logarithm,
zero at $z=0$, therefore has limiting phase equal to the indicated
continuous gamma phase, rather than that phase plus a nonzero
multiple of $2\pi$. Equivalently that possible integer is locally
constant in $\lambda$ and is zero at $\lambda=0$.

The pair of pants has its reflection fixing the perpendicular
seams. After capping, this antiholomorphic involution fixes
$0,1,\infty$, so it is complex conjugation. The perpendicular
seam from the cusp at $1$ to the geodesic boundary is the terminal
segment of the real interval $(0,1)$. Its developing image lies
on the unit circle, by (AQ.67.4). The cusp endpoint has phase
$\vartheta$. By (AQ.67.5), approaching that cusp from the core
uses phases less than $\vartheta$. Hence the developing half-plane
of the core is the half-plane with argument interval
$(\vartheta-\pi,\vartheta)$. The axis of the boundary monodromy,
whose fixed points are $0,\infty$, is its ray of argument
$\vartheta-\pi/2$. The seam meets that axis perpendicularly.
Starting from the cusp, this is the first intersection of the
seam with the axis; the further formal intersections obtained
after crossing the boundary line of that half-plane are not in
the hyperbolic core. This selects the lifted phase
$\vartheta-\pi/2$, with no integer ambiguity.

Choose $w=1$ at this perpendicular foot. On a collar of the
geodesic, and thence by analytic continuation into its cap,
\[
 w=\exp\left(\frac{\log D-i(\vartheta-\pi/2)}{i\lambda}\right).
\]
Here the logarithm has the lifted phase just specified.
Its boundary modulus is one and its boundary argument has the
normalized arclength, so this is precisely the prescribed cap
coordinate. At zero, (AQ.67.4) gives
\[
 w=\exp\left(-\frac{\vartheta-\pi/2}{\lambda}\right)z+O(z^2),
 \qquad
 r(\ell)=\exp\left(\frac{\vartheta-\pi/2}{\lambda}\right).
                                                               \tag{AQ.67.7}
\]
The duplication formula, DLMF (5.5.5), yields
$\vartheta=2\lambda\log2+2\phi(\lambda/2)$, proving
(AQ.67.1).

It remains to establish the exact range and the nonvanishing
derivative, rather than only an upper bound. The digamma integral
DLMF (5.9.16) gives, for $x\geq0$,
\[
 \phi'(x)=\operatorname{Re}\bigl(\psi(1+ix)
                            -\psi(\tfrac12+ix)\bigr)
           =\int_0^\infty\frac{\cos(xt)}{e^{t/2}+1}\,dt.
                                                               \tag{AQ.67.8}
\]
Let $f(t)=(e^{t/2}+1)^{-1}$. Then
\[
 f''(t)=\frac{e^{t/2}(e^{t/2}-1)}{4(e^{t/2}+1)^3}>0
               \quad(t>0).
\]
The function and its derivatives decay exponentially. Two
integrations by parts, with their endpoint terms included, give
\[
 \phi'(x)=\frac1{x^2}\int_0^\infty f''(t)(1-\cos xt)\,dt>0
                  \quad(x>0).                                \tag{AQ.67.9}
\]
Stirling's expansion for the same analytic logarithms, DLMF
(5.11.1), or the corresponding gamma-ratio expansion (5.11.13),
gives $\phi(x)\longrightarrow\pi/4$ as $x\to\infty$.
Since $\phi(0)=0$, it follows that
$0<\phi(x)<\pi/4$ for $x>0$. Therefore
\[
 \frac{d}{dx}\log r(4\pi x)
       =\frac{x\phi'(x)+\pi/4-\phi(x)}{x^2}>0.                 \tag{AQ.67.10}
\]
At zero, $\phi'(0)=2\log2$, so
\[
 r(\ell)=16\exp\left(-\frac{\pi^2}{\ell}+O(\ell^2)\right)
              \quad(\ell\downarrow0),
 \qquad r(\ell)\longrightarrow4\quad(\ell\to\infty).
\]
Real analyticity and (AQ.67.10) prove that $r$ is the asserted
diffeomorphism.

Finally Fenchel--Nielsen gluing of the two labelled pairs of
pants has exactly one length $\ell>0$ and one twist modulo
$\ell$, equivalently an angle $\theta$ modulo $2\pi$.
The two derivative vectors each have norm $r(\ell)$, and
the twist contributes their tensor's unit phase. This proves
(AQ.67.2). Its inverse is
$\ell=r^{-1}(\sqrt{|q|})$, $\theta=\arg q$ with the latter
viewed as a circle-valued coordinate. In logarithmic polar
coordinates the real Jacobian is diagonal, with entries
$2r'(\ell)/r(\ell)>0$ and $1$. Both the punctured disk and
$\mathbb C^\times$ deformation retract to the same circle
$|q|=1$; the inclusion therefore induces an isomorphism on
$\pi_1$, taking one full twist to the positive generator.
\end{proof}
